% ============================================================================
% MÓDULO 3 — SIMULADO PREPARATÓRIO PARA A AVALIAÇÃO TEÓRICA (SEMANA 10)
% Abrange integralmente os Módulos 1 a 3 (Capítulos 1 a 9)
% CHF662 — ANÁLISE AMBIENTAL (DCHF/UEFS)
% ============================================================================

\clearpage
\phantomsection
\addcontentsline{toc}{chapter}{Simulado Preparatório para a Avaliação Teórica (Semana 10: Módulos 1 a 3)}
\chapter*{Simulado Preparatório para a Avaliação Teórica: Módulos 1 a 3}
\markboth{SIMULADO, SEMANA 10 (MÓDULOS 1 A 3)}{SIMULADO, SEMANA 10 (MÓDULOS 1 A 3)}

\noindent\tagEmenta\ \tagPratica

\vspace{6pt}

\begin{boxConceito}[Instruções de Aplicação e Padrão da Avaliação Teórica (Semana 10)]
  Este simulado foi estruturado exatamente no modelo da \textbf{Avaliação Teórica Presencial da Semana 10 ($3{,}0$ pontos)}, integrando todos os conteúdos dos \textbf{Módulos 1, 2 e 3 (Capítulos 1 a 9)}: crise ambiental e epistemologia (Leff; Porto-Gonçalves; Mendonça), conceitos e AIA (Sánchez), sistemas abertos e escalas, bacias hidrográficas e morfometria, morfopedologia e declividade, geossistemas (Bertrand), ecodinâmica (Tricart) e fragilidade ambiental (Ross).
  \begin{itemize}[leftmargin=1.5em, itemsep=2pt]
    \item \textbf{Tempo sugerido:} $150$ minutos ($50$ min para as $10$ questões objetivas + $100$ min para as $4$ questões dissertativas integradoras).
    \item \textbf{Pontuação simulada (referência da Prova Presencial = $3{,}0$ pontos):} na Parte II, cada uma das $4$ questões dissertativas vale $0{,}75$ ponto ($4 \times 0{,}75 = 3{,}00$ pontos), avaliadas pelo Espelho Dissertativo ao final. As $10$ questões objetivas da Parte I funcionam como treino de fixação e diagnóstico conceitual estilo ENADE e concursos.
    \item \textbf{Regra de ouro:} resolva integralmente sem consulta antes de conferir o Gabarito Comentado e preencher a \textbf{Tabela de Diagnóstico de Revisão}.
  \end{itemize}
\end{boxConceito}

% ============================================================================
% PARTE I — 10 QUESTÕES OBJETIVAS A–E (ESTILO ENADE E CONCURSOS AMBIENTAIS)
% ============================================================================
\section*{Parte I: Questões Objetivas (Estilo ENADE e Concursos de Geografia e Análise Ambiental)}

\begin{enumerate}[leftmargin=1.8em, itemsep=6pt]
  \item \textbf{[Autoral, Estilo ENADE / Epistemologia Ambiental (Cap. 1)]} Ao analisar as raízes da crise ecológica contemporânea, Enrique Leff (2001) argumenta que a degradação dos ecossistemas e a exclusão social não resultam apenas de falhas técnicas pontuais. Segundo esse autor, a crise ambiental expressa fundamentalmente:
  \begin{alternativas}
    \item Um fenômeno natural cíclico desvinculado dos modos de produção e consumo.
    \item Uma crise civilizatória e de racionalidade científica e econômica, marcada pela fragmentação do saber e pela mercantilização ilimitada da natureza.
    \item Uma limitação exclusiva da engenharia hidráulica em regiões semiáridas.
    \item Um problema estritamente demográfico solucionável sem alteração do modelo produtivo.
    \item Uma consequência inevitável da conservação de unidades de proteção integral.
  \end{alternativas}
  \espacoResolucao[1.1cm]

  \item \textbf{[Autoral, Estilo ENADE / Sociedade--Natureza e Geografia Socioambiental (Caps. 2 e 3)]} Articulando as contribuições de Porto-Gonçalves (2006) e Mendonça (2001) sobre a relação sociedade--natureza na Geografia, assinale a alternativa correta:
  \begin{alternativas}
    \item A Geografia Socioambiental propõe separar rigidamente a Geografia Física da Geografia Humana para preservar a neutralidade científica.
    \item O conceito de natureza é uma construção histórica e cultural; a separação moderna sujeito--objeto legitimou a apropriação desigual dos recursos e os conflitos socioambientais territoriais.
    \item Os conflitos socioambientais ocorrem apenas em áreas metropolitanas industriais, inexistindo no espaço agrário ou semiárido.
    \item A abordagem socioambiental nega a materialidade física do relevo, dos solos e da água.
    \item O determinismo ambiental explica integralmente os conflitos por água no Semiárido brasileiro.
  \end{alternativas}
  \espacoResolucao[1.1cm]

  \item \textbf{[Autoral, Estilo Concurso Analista Ambiental / Conceitos Fundamentais (Cap. 4)]} Em uma comunidade ribeirinha assentada sobre planície de inundação assoreada, a ocorrência de cheias sazonais representa um perigo físico elevado, agravado pela precariedade das habitações e pela ausência de drenagem urbana. À luz de Sánchez (2013) e Tauk-Tornisielo (1996), é correto afirmar que:
  \begin{alternativas}
    \item Perigo e risco são sinônimos perfeitos, pois ambos existem mesmo sem população ou bens expostos.
    \item A resiliência ecológica aumenta quando o canal fluvial é impermeabilizado e retificado.
    \item O risco socioambiental resulta da combinação entre a ameaça (perigo de cheia) e a vulnerabilidade das pessoas e infraestruturas expostas.
    \item Degradação ambiental refere-se exclusivamente a processos naturais anteriores à ocupação humana.
    \item Impacto ambiental inclui apenas alterações negativas, nunca efeitos positivos.
  \end{alternativas}
  \espacoResolucao[1.1cm]

  \item \textbf{[Autoral, Estilo Concurso IBAMA / AIA e Matriz de Leopold Ponderada (Cap. 5)]} Na avaliação de impacto ambiental de uma rodovia rural, uma coluna da Matriz de Leopold Ponderada ($P_{ij} = s_{ij} \cdot M_{ij} \cdot I_{ij}$) apresentou três interações: supressão de mata ciliar $(-, M=8, I=7)$, erosão em cortes de talude $(-, M=4, I=5)$ e melhoria do escoamento da produção local $(+, M=5, I=6)$. Os valores de $\Sigma^{-}$, $\Sigma^{+}$ e do total da coluna $T$ são, respectivamente:
  \begin{alternativas}
    \item $\Sigma^{-} = -76$, $\Sigma^{+} = +30$, $T = -46$.
    \item $\Sigma^{-} = -56$, $\Sigma^{+} = +30$, $T = -26$.
    \item $\Sigma^{-} = -76$, $\Sigma^{+} = +11$, $T = -65$.
    \item $\Sigma^{-} = -24$, $\Sigma^{+} = +30$, $T = +6$.
    \item $\Sigma^{-} = -76$, $\Sigma^{+} = +30$, $T = -106$.
  \end{alternativas}
  \espacoResolucao[1.3cm]

  \item \textbf{[Autoral, Estilo Concurso Geógrafo / Sistemas Abertos e Falácia de Escala (Cap. 6)]} Uma microbacia hidrográfica de $800\,\text{ha}$ apresenta erosão bruta estimada em vertentes de $15{,}0\,\text{t/(ha}\cdot\text{ano)}$, mas a estação hidrossedimentológica no exutório mede uma descarga sólida de $3\,000\,\text{t/ano}$. Sobre esse balanço e o conceito de falácia de escala, assinale a alternativa correta:
  \begin{alternativas}
    \item A erosão bruta total é de $3\,000\,\text{t/ano}$ e não há retenção interna de sedimentos.
    \item A Razão de Entrega de Sedimentos é $\text{SDR} = 0{,}75$, pois $75\%$ do solo erodido alcança o exutório.
    \item A Razão de Entrega de Sedimentos é $\text{SDR} = 0{,}25$, indicando que $9\,000\,\text{t/ano}$ ficam retidas em sopés e planícies; assumir que toda a erosão de parcela chega ao exutório configura falácia de escala.
    \item O resultado indica erro de medição, pois $\text{SDR}$ em sistemas abertos é sempre igual a $1{,}0$.
    \item A retenção interna de sedimentos equivale a $3\,000\,\text{t/ano}$.
  \end{alternativas}
  \espacoResolucao[1.3cm]

  \item \textbf{[Autoral, Estilo Concurso Perito Ambiental / Morfometria de Bacias (Cap. 7)]} Uma bacia hidrográfica apresenta área $A = 25\,\text{km}^2$, perímetro $P = 22\,\text{km}$ e comprimento total da rede de drenagem $L_t = 40\,\text{km}$. Utilizando $D_d = L_t / A$ e $K_c = 0{,}28 \cdot P / \sqrt{A}$, os valores corretos da Densidade de Drenagem ($D_d$) e do Coeficiente de Compacidade ($K_c$) são, respectivamente:
  \begin{alternativas}
    \item $D_d = 0{,}625\,\text{km/km}^2$ e $K_c = 1{,}000$.
    \item $D_d = 1{,}60\,\text{km/km}^2$ e $K_c = 1{,}232$.
    \item $D_d = 1{,}60\,\text{km/km}^2$ e $K_c = 0{,}880$.
    \item $D_d = 2{,}20\,\text{km/km}^2$ e $K_c = 1{,}232$.
    \item $D_d = 1{,}82\,\text{km/km}^2$ e $K_c = 2{,}464$.
  \end{alternativas}
  \espacoResolucao[1.3cm]

  \item \textbf{[Autoral, Estilo ENADE / Bacias, Áreas Ripárias e Conectividade (Cap. 7)]} Em uma microbacia rural do Semiárido baiano, a substituição da vegetação ripária nativa por pastagem compactada e abertura de estradas não pavimentadas cruzando talvegues provoca, de forma integrada:
  \begin{alternativas}
    \item Aumento da infiltração profunda, redução dos picos de cheia e estabilização das margens.
    \item Desconexão total entre vertentes e canais, eliminando o transporte de sedimentos.
    \item Aumento da conectividade hidrossedimentológica, elevação dos picos de cheia, erosão de margens e assoreamento acelerado dos açudes a jusante.
    \item Redução imediata da energia do escoamento superficial durante chuvas convectivas.
    \item Transformação de canais de 1ª ordem em rios perenes de 5ª ordem pelo critério de Strahler.
  \end{alternativas}
  \espacoResolucao[1.1cm]

  \item \textbf{[Autoral, Estilo Concurso Geografia / Morfopedologia e Declividade (Cap. 8)]} Em uma carta topográfica na escala $1:10\,000$, a distância medida entre duas curvas de nível com desnível $\Delta h = 45\,\text{m}$ é de $d_{\text{mapa}} = 1{,}5\,\text{cm}$. Sabendo que a encosta possui solo raso sobre rocha pouco permeável, a declividade ($S\%$ e $\theta$) e a classe de declividade segundo Ross (1994) são:
  \begin{alternativas}
    \item $S = 30{,}0\%$; $\theta \approx 16{,}70^\circ$; limite superior da Classe 4 de Ross (Forte, $20\text{--}30\%$), com risco de escorregamento translacional.
    \item $S = 3{,}0\%$; $\theta \approx 1{,}72^\circ$; Classe 1 de Ross (Muito Fraca).
    \item $S = 15{,}0\%$; $\theta \approx 8{,}53^\circ$; Classe 3 de Ross (Média).
    \item $S = 45{,}0\%$; $\theta \approx 24{,}23^\circ$; Classe 5 de Ross (Muito Forte).
    \item $S = 66{,}7\%$; $\theta \approx 33{,}69^\circ$; Classe 5 de Ross (Muito Forte).
  \end{alternativas}
  \espacoResolucao[1.3cm]

  \item \textbf{[Autoral, Estilo ENADE / Geossistemas e Ecodinâmica (Cap. 9)]} Articulando a taxonomia da paisagem de Bertrand (2004) e a Ecodinâmica de Tricart (1977), analise as afirmativas abaixo:
  \begin{enumerate}[label=\Roman*., leftmargin=2.0em, itemsep=2pt]
    \item O geossistema resulta da combinação dinâmica entre potencial ecológico, exploração biológica e ação antrópica em unidade territorial concreta.
    \item Nos meios estáveis de Tricart, a pedogênese supera amplamente a morfogênese, favorecendo solos profundos sob cobertura vegetal conservada.
    \item Nos meios fortemente instáveis, o balanço morfogênese--pedogênese pende fortemente para a morfogênese, com remoção ativa de solo e ocorrência de ravinas ou movimentos de massa.
  \end{enumerate}
  É correto o que se afirma em:
  \begin{alternativas}
    \item I, apenas.
    \item II, apenas.
    \item I e II, apenas.
    \item II e III, apenas.
    \item I, II e III.
  \end{alternativas}
  \espacoResolucao[1.1cm]

  \item \textbf{[Autoral, Estilo Concurso Órgão Ambiental / Fragilidade Ambiental de Ross (Cap. 9)]} Na Análise Empírica da Fragilidade Ambiental de Ross (1994), uma encosta possui declividade $S = 26\%$ ($D_R = 4$), Planossolo ($D_S = 4$), chuvas sazonais intensas ($D_C = 4$) e cobertura por floresta secundária avançada ($D_U = 2$). Os códigos e as classes síntese de Fragilidade Potencial ($FP$) e Fragilidade Emergente ($FE$) dessa unidade são:
  \begin{alternativas}
    \item $FP = 444$ (Classe 4, Forte) e $FE = 4442$ (Classe 3, Média).
    \item $FP = 444$ (Classe 4, Forte) e $FE = 4442$ (Classe 5, Muito Forte).
    \item $FP = 222$ (Classe 2, Fraca) e $FE = 2224$ (Classe 4, Forte).
    \item $FP = 442$ (Classe 3, Média) e $FE = 4422$ (Classe 2, Fraca).
    \item $FP = 554$ (Classe 5, Muito Forte) e $FE = 5542$ (Classe 4, Forte).
  \end{alternativas}
  \espacoResolucao[1.3cm]
\end{enumerate}

% ============================================================================
% PARTE II — 4 QUESTÕES DISSERTATIVAS INTEGRADORAS (MODELO PROVA CHF662)
% ============================================================================
\section*{Parte II: Questões Dissertativas Integradoras (Modelo da Avaliação Teórica CHF662: $3{,}0$ Pontos)}

\begin{enumerate}[leftmargin=1.8em, itemsep=6pt, resume]
  \item \textbf{[Questão Dissertativa 1, Epistemologia, Sociedade--Natureza e Conceitos de Análise Ambiental ($0{,}75$ ponto)]}
  Em um município do Semiárido baiano, a expansão desordenada de empreendimentos energéticos e minerários sobre terras de uso tradicional provocou desmatamento de Caatinga, assoreamento de aguadas comunitárias e perda de acesso das famílias à água.
  \begin{enumerate}[label=(\alph*), leftmargin=1.8em, itemsep=2pt]
    \item À luz de \textbf{Leff (2001)} e \textbf{Porto-Gonçalves (2006)}, explique por que esse cenário expressa uma crise de racionalidade e um conflito socioambiental, e não apenas um problema técnico isolado. ($0{,}25$ ponto)
    \item Com base em \textbf{Mendonça (2001)}, justifique por que a Geografia Socioambiental exige articular os processos físico-naturais e os processos sociais nesse diagnóstico. ($0{,}25$ ponto)
    \item Diferencie rigorosamente, segundo \textbf{Sánchez (2013)}, os conceitos de \textbf{impacto ambiental}, \textbf{degradação ambiental}, \textbf{vulnerabilidade} e \textbf{resiliência} aplicados ao caso descrito. ($0{,}25$ ponto)
  \end{enumerate}
  \espacoResolucao[3.6cm]

  \item \textbf{[Questão Dissertativa 2, AIA, Sistemas Abertos, Morfometria e Conectividade em Bacias ($0{,}75$ ponto)]}
  Uma microbacia hidrográfica rural possui área $A = 16\,\text{km}^2$, perímetro $P = 20\,\text{km}$ e comprimento total de canais de drenagem $L_t = 36\,\text{km}$. Prevê-se a implantação de um distrito agroindustrial na média encosta sem bacias de contenção nem preservação de APPs ripárias.
  \begin{enumerate}[label=(\alph*), leftmargin=1.8em, itemsep=2pt]
    \item Calcule a \textbf{Densidade de Drenagem ($D_d$)} e o \textbf{Coeficiente de Compacidade ($K_c$)}, interpretando a eficiência de escoamento e a forma da bacia quanto à susceptibilidade a cheias. ($0{,}25$ ponto)
    \item Explique como a remoção da mata ciliar e a impermeabilização alteram a \textbf{conectividade hidrossedimentológica} e o balanço entradas--saídas da bacia como sistema aberto. ($0{,}25$ ponto)
    \item Indique duas ações concretas aplicando os dois primeiros degraus da \textbf{Hierarquia de Mitigação} (Evitar e Minimizar) no processo de AIA desse empreendimento. ($0{,}25$ ponto)
  \end{enumerate}
  \espacoResolucao[3.8cm]

  \item \textbf{[Questão Dissertativa 3, Morfopedologia, Cálculo de Declividade e Movimentos de Massa ($0{,}75$ ponto)]}
  Durante vistoria em uma encosta urbana ocupada sem projeto geotécnico, mediu-se com trena sobre a rampa uma distância inclinada $d_{\text{inc}} = 130\,\text{m}$ entre a base e o topo, com desnível vertical $\Delta h = 50\,\text{m}$. O solo local é um manto argilo-arenoso espesso que recebe lançamento contínuo de águas servidas.
  \begin{enumerate}[label=(\alph*), leftmargin=1.8em, itemsep=2pt]
    \item Calcule a \textbf{distância horizontal real ($d_h$)}, a \textbf{declividade em porcentagem ($S\%$)} e o \textbf{ângulo em graus ($\theta$)}, classificando a encosta segundo a \textbf{EMBRAPA} e \textbf{Ross (1994)}. ($0{,}30$ ponto)
    \item Explique por que utilizar $d_{\text{inc}} = 130\,\text{m}$ diretamente no denominador de $S(\%)$ configuraria erro metodológico. ($0{,}15$ ponto)
    \item Identifique o tipo de \textbf{movimento de massa} mais provável nesse perfil espesso e explique o papel da \textbf{poropressão ($u$)} causada pelas águas servidas na ruptura do talude. ($0{,}30$ ponto)
  \end{enumerate}
  \espacoResolucao[3.8cm]

  \item \textbf{[Questão Dissertativa 4, Geossistemas, Ecodinâmica e Fragilidade Ambiental de Ross ($0{,}75$ ponto)]}
  Uma unidade geoambiental de meia encosta no Semiárido baiano apresenta declividade $S = 34\%$, Argissolo com mudança textural abrupta ($D_S = 4$), chuvas sazonais torrenciais ($D_C = 4$) e encontra-se atualmente com solo exposto após desmatamento ($D_U = 5$).
  \begin{enumerate}[label=(\alph*), leftmargin=1.8em, itemsep=2pt]
    \item Enquadre essa vertente na escala taxonômica de \textbf{Bertrand (2004)} (Geossistema vs. Geofácies) e classifique seu estado atual pela \textbf{Ecodinâmica de Tricart (1977)}, justificando o balanço morfogênese--pedogênese. ($0{,}25$ ponto)
    \item Calcule o código numérico e a classe de \textbf{Fragilidade Potencial ($FP$)} e o código numérico e a classe de \textbf{Fragilidade Emergente ($FE$)} segundo \textbf{Ross (1994)}. ($0{,}30$ ponto)
    \item Caso um projeto de recuperação recomponha a vegetação até atingir cobertura arbustivo-arbórea densa ($D_U = 2$), qual passará a ser o novo código e a nova classe de $FE$? Explique o significado ambiental dessa mudança. ($0{,}20$ ponto)
  \end{enumerate}
  \espacoResolucao[3.8cm]
\end{enumerate}

% ============================================================================
% GABARITO COMENTADO, ESPELHO DISSERTATIVO E DIAGNÓSTICO DE REVISÃO
% ============================================================================
\clearpage
\section*{Gabarito Comentado: Simulado Preparatório (Semana 10)}

\begin{center}
\renewcommand{\arraystretch}{1.2}
\small
\begin{tabularx}{\textwidth}{@{}c c c c c c c c c c@{}}
  \toprule
  \rowcolor{uefsNavy}
  \textcolor{white}{\bfseries Q1} & \textcolor{white}{\bfseries Q2} & \textcolor{white}{\bfseries Q3} & \textcolor{white}{\bfseries Q4} & \textcolor{white}{\bfseries Q5} & \textcolor{white}{\bfseries Q6} & \textcolor{white}{\bfseries Q7} & \textcolor{white}{\bfseries Q8} & \textcolor{white}{\bfseries Q9} & \textcolor{white}{\bfseries Q10} \\
  \midrule
  \textbf{B} & \textbf{B} & \textbf{C} & \textbf{A} & \textbf{C} & \textbf{B} & \textbf{C} & \textbf{A} & \textbf{E} & \textbf{A} \\
  \bottomrule
\end{tabularx}
\end{center}

\begin{boxGabarito}[Comentário Detalhado das Questões Objetivas (Q1 a Q10)]
  \begin{itemize}[leftmargin=1.5em, itemsep=3pt]
    \item \textbf{Q1, (B):} para Leff (2001), a crise ambiental é uma crise civilizatória e de racionalidade científica/econômica que fragmentou o saber e tratou a natureza como objeto mercantil infinito.
    \item \textbf{Q2, (B):} Porto-Gonçalves (2006) evidencia a historicidade da ideia de natureza e a gênese dos conflitos territoriais; Mendonça (2001) integra o quadro físico e a dinâmica social na Geografia Socioambiental.
    \item \textbf{Q3, (C):} perigo é a ameaça física (cheia); risco existe quando há população/bens vulneráveis expostos a esse perigo ($R = P \times V$).
    \item \textbf{Q4, (A):} células: $-(8 \times 7) = -56$, $-(4 \times 5) = -20$ e $+(5 \times 6) = +30$. Portanto, $\Sigma^{-} = -56 - 20 = -76$, $\Sigma^{+} = +30$ e $T = -76 + 30 = -46$.
    \item \textbf{Q5, (C):} erosão bruta total $= 15{,}0 \times 800 = 12\,000\,\text{t/ano}$. Como $3\,000\,\text{t/ano}$ saem no exutório, $\text{SDR} = 3\,000 / 12\,000 = \mathbf{0{,}25}$ ($25\%$), com $9\,000\,\text{t/ano}$ retidas na bacia.
    \item \textbf{Q6, (B):} $D_d = 40 / 25 = \mathbf{1{,}60\,\text{km/km}^2}$; $K_c = 0{,}28 \times 22 / \sqrt{25} = 6{,}16 / 5 = \mathbf{1{,}232}$.
    \item \textbf{Q7, (C):} estradas e remoção da mata ripária conectam diretamente as vertentes aos canais, acelerando o escoamento, elevando picos de cheia e assoreando reservatórios a jusante.
    \item \textbf{Q8, (A):} $d_h = 1{,}5\,\text{cm} \times 10\,000 = 15\,000\,\text{cm} = 150\,\text{m}$. Logo $S = (45 / 150) \times 100 = \mathbf{30{,}0\%}$ e $\theta = \arctan(0{,}30) \approx \mathbf{16{,}70^\circ}$ (limite superior da Classe 4 de Ross, $20\text{--}30\%$). Em solo raso sobre rocha, o mecanismo crítico é o escorregamento translacional.
    \item \textbf{Q9, (E):} as três afirmativas (I, II e III) sintetizam fielmente os conceitos de geossistema (Bertrand) e os meios estáveis e fortemente instáveis (Tricart).
    \item \textbf{Q10, (A):} $D_R=4, D_S=4, D_C=4 \implies FP = \mathbf{444}$, média $4{,}0$ (Classe 4, Forte). Cruzando $FP=4$ com $D_U=2$ na Matriz de Ross, obtém-se $FE = \mathbf{4442}$ (\textbf{Classe 3, Média}).
  \end{itemize}
\end{boxGabarito}

\section*{Espelho Dissertativo com Critérios de Pontuação (Q11 a Q14: Total: $3{,}0$ Pontos)}

\begin{boxGabarito}[Espelho: Questão 11 ($0{,}75$ ponto) e Questão 12 ($0{,}75$ ponto)]
  \textbf{Questão 11 ($0{,}75$ ponto):}
  \begin{itemize}[leftmargin=1.5em, itemsep=2.5pt]
    \item \textbf{(a) $0{,}25$ pt:} explicar que a crise decorre da racionalidade econômica instrumental que reduz a Caatinga e a água a mercadorias (Leff, 2001) e gera disputa territorial/simbólica pelo acesso aos bens comuns (Porto-Gonçalves, 2006).
    \item \textbf{(b) $0{,}25$ pt:} demonstrar que a Geografia Socioambiental (Mendonça, 2001) supera a dicotomia físico vs. humano ao integrar a dinâmica hidrológica/erosiva às relações de poder e uso do território.
    \item \textbf{(c) $0{,}25$ pt:} definir \emph{impacto ambiental} (alteração mensurável positiva ou negativa da qualidade ambiental), \emph{degradação} (perda adversa de funcionalidade ecológica/produtiva), \emph{vulnerabilidade} (grau de suscetibilidade social/ecológica ao dano) e \emph{resiliência} (capacidade de absorver perturbações e recuperar funções).
  \end{itemize}
  \vspace{4pt}
  \textbf{Questão 12 ($0{,}75$ ponto):}
  \begin{itemize}[leftmargin=1.5em, itemsep=2.5pt]
    \item \textbf{(a) $0{,}25$ pt:} $D_d = 36 / 16 = \mathbf{2{,}25\,\text{km/km}^2}$ (drenagem moderada a alta); $K_c = 0{,}28 \times 20 / \sqrt{16} = 5{,}60 / 4 = \mathbf{1{,}400}$ (forma moderadamente alongada, menos propensa a picos extremos que uma bacia circular, embora exija controle de drenagem).
    \item \textbf{(b) $0{,}25$ pt:} a impermeabilização reduz a infiltração e acelera o escoamento superficial; sem a mata ciliar, aumenta a conectividade entre vertente e canal, elevando o pico de vazão e a carga sólida exportada no exutório.
    \item \textbf{(c) $0{,}25$ pt:} \emph{Evitar:} realocar a poligonal do distrito para fora das APPs ripárias e setores mais declivosos; \emph{Minimizar:} implantar bacias de retenção/decantação, drenagem perimetral e pavimentos permeáveis.
  \end{itemize}
\end{boxGabarito}

\begin{boxGabarito}[Espelho: Questão 13 ($0{,}75$ ponto) e Questão 14 ($0{,}75$ ponto)]
  \textbf{Questão 13 ($0{,}75$ ponto):}
  \begin{itemize}[leftmargin=1.5em, itemsep=2.5pt]
    \item \textbf{(a) $0{,}30$ pt:} $d_h = \sqrt{130^2 - 50^2} = \sqrt{16\,900 - 2\,500} = \sqrt{14\,400} = \mathbf{120{,}0\,\text{m}}$. Declividade: $S = (50 / 120) \times 100 \approx \mathbf{41{,}67\%}$; $\theta = \arctan(50 / 120) \approx \mathbf{22{,}62^\circ}$. Classificação: \textbf{EMBRAPA Forte ondulado ($20\text{--}45\%$)} e \textbf{Ross Classe 5, Muito Forte ($>30\%$)}.
    \item \textbf{(b) $0{,}15$ pt:} $d_{\text{inc}}$ é a hipotenusa do triângulo retângulo; usá-la no denominador calcula $\sin\theta$ em vez de $\tan\theta$, subestimando a declividade real ($38{,}46\%$ em vez de $41{,}67\%$).
    \item \textbf{(c) $0{,}30$ pt:} em solo argilo-arenoso espesso, o mecanismo provável é o \textbf{escorregamento rotacional} (superfície de ruptura curva). As águas servidas saturam os poros e elevam a \textbf{poropressão ($u$)}, reduzindo a tensão efetiva e o atrito interno entre os grãos, facilitando o colapso.
  \end{itemize}
  \vspace{4pt}
  \textbf{Questão 14 ($0{,}75$ ponto):}
  \begin{itemize}[leftmargin=1.5em, itemsep=2.5pt]
    \item \textbf{(a) $0{,}25$ pt:} a vertente homogênea corresponde ao \textbf{Geofácies} (inserido em um \textbf{Geossistema} de sub-bacia/escala $1:50.000$). Pela Ecodinâmica de Tricart, encontra-se em \textbf{Meio Fortemente Instável}, pois a alta declividade e o solo exposto fazem a \textbf{morfogênese superar amplamente a pedogênese} ($M \gg P$).
    \item \textbf{(b) $0{,}30$ pt:} $D_R = 5$ ($34\% > 30\%$), $D_S = 4$, $D_C = 4$, $D_U = 5$. Código $FP = \mathbf{544}$; média $I_{FP} = (5+4+4)/3 = 13/3 \approx 4{,}33 \implies$ \textbf{Classe $FP = 4$ (Forte)}. Código $FE = \mathbf{5445}$; cruzando $FP=4 \times D_U=5$ na Matriz de Ross $\implies$ \textbf{Classe $FE = 5$ (Muito Forte)}.
    \item \textbf{(c) $0{,}20$ pt:} com $D_U = 2$, o novo código será $FE = \mathbf{5442}$. Cruzando $FP=4 \times D_U=2$ na Matriz de Ross, a classe cai de $5$ (Muito Forte) para \textbf{Classe $FE = 3$ (Média)}. A recomposição vegetal amortece a energia das gotas, amplia a infiltração e reduz a fragilidade real, embora a fragilidade potencial natural ($FP=4$) exija conservação permanente.
  \end{itemize}
\end{boxGabarito}

\section*{Tabela de Diagnóstico de Revisão: Preparação Final para a Semana 10}

\begin{center}
\renewcommand{\arraystretch}{1.22}
\small
\begin{tabularx}{\textwidth}{@{}c c p{4.3cm} >{\raggedright\arraybackslash}X@{}}
  \toprule
  \rowcolor{uefsNavy}
  \textcolor{white}{\bfseries Questão} & \textcolor{white}{\bfseries Módulo / Cap.} & \textcolor{white}{\bfseries Conteúdo Avaliado} & \textcolor{white}{\bfseries Ação Corretiva de Revisão} \\
  \midrule
  Q1, Q2, Q11(a--b) & Mód. 1 (Caps. 1--3) & Crise ambiental, sociedade--natureza e Geografia Socioambiental & Reler sínteses de Leff, Porto-Gonçalves e Mendonça nos Caps. 1 a 3. \\
  \rowcolor{bgGray}
  Q3, Q11(c) & Mód. 1 (Cap. 4) & Impacto, degradação, risco, vulnerabilidade e resiliência & Revisar quadros conceituais e erros comuns do Cap. 4. \\
  Q4, Q5, Q12(c) & Mód. 2 (Caps. 5--6) & Matriz de Leopold, Hierarquia de Mitigação, Sistemas e SDR & Refazer os cálculos de Leopold Ponderada, SDR e Escalas (Caps. 5 e 6). \\
  \rowcolor{bgGray}
  Q6, Q7, Q12(a--b) & Mód. 3 (Cap. 7) & Morfometria ($D_d$, $K_c$), Strahler e conectividade em bacias & Revisar quadro de equações e Exemplos 7.2 a 7.5 do Cap. 7. \\
  Q8, Q13 & Mód. 3 (Cap. 8) & Topossequência, declividade ($S\%$, $\theta$), trena e movimentos de massa & Refazer Exemplos 8.2, 8.3 e 8.5 e Erros Comuns 1--3 do Cap. 8. \\
  \rowcolor{bgGray}
  Q9, Q10, Q14 & Mód. 3 (Cap. 9) & Bertrand, Tricart (Ecodinâmica) e Ross ($FP$ e $FE$) & Revisar Tabelas 9.2--9.4, Figura~\ref{fig:ecodinamica_ross} e Exemplos 9.3--9.5. \\
  \bottomrule
\end{tabularx}
\end{center}
